\subsection{The first nonfree bar repair extends to prescribed targets}

\begin{theorem}[D-030: \status{Proved}; extension through bar arity two]
\label{sa08:D030}
Let \(P\) be a simplicial \(k\)-module with projective
normalized terms, and let
\(x:P\to A^{w,1}\) be a prescribed simplicial linear map.
Assume that \(A^{w,1}\) is degreewise flat and horizontally
acyclic. In weight \(2w\), assume that the vertical
augmentation \(\underline{k}\to A^{2w,*}\) is an
isomorphism on \(H^0\), that \(H^j(A^{2w,*})=0\) for
\(1\le j\le4\), and that all rows \(A^{2w,j}\),
\(0\le j\le4\), are horizontally acyclic.

The map of the original odd-to-even disk, given by
\(x\) and \(y=dx\), extends to the bar closure
\(U_1(P)\) through parameter arity two.
Precisely, there are maps from its arities zero, one
and two to \(A^{0,*},A^{w,*},A^{2w,*}\), compatible
with the differential and with every product and
divided power whose output arity is at most two.
The arity-one map is unchanged. All new parameters
in this first repair have form degree at least two.

This is a statement about a partial map through arity
two. It asserts neither an extension over every higher
arity nor an extension relative to separately prescribed
boundary repairs.
\end{theorem}

\begin{proof}
Work in the abelian category
\(\mathcal S_k=s\operatorname{Mod}_k\).
The quadratic part of the original free disk is
\[
 C(P)=\bigl(\Lambda^2P\longrightarrow P\otimes P
                         \longrightarrow\Gamma^2P\bigr)
\]
in form degrees two, three and four. Its prescribed
map to \(A^{2w,*}\) is the map of products and divided
squares of \(x,dx\).
Statement~\ref{sa08:C23}, with \(q=1,m=2\), applies
under exactly the stated finite-range hypotheses.
It gives a strict simplicial linear cochain homotopy
\[
 h^i:C^i(P)\longrightarrow A^{2w,i-1},\qquad
                  dh+hd=f.
 \tag{D-30a}
\]

The arity-two part \(U\) of the bar closure has terms
\[
 U^2=P\otimes P,\qquad
 U^3=(P\otimes P)\oplus\Gamma^2P,\qquad
 U^4=\Gamma^2P.
 \tag{D-30b}
\]
The old disk embeds by the alternating-tensor map
in degree two, by the coefficient-times-bar summand
in degree three, and by the identity in degree four.
These maps are injective for projective modules:
the free case follows on basis tensors, and the
projective case follows by taking retracts.
The quotient complex \(R=U/C(P)\) has
\[
 R^2=\operatorname{Sym}^2P,
 \qquad R^3=\Gamma^2P,
 \qquad R^i=0\quad(i\ne2,3).
 \tag{D-30c}
\]
Its differential is, up to the bar sign, the norm
map \(\operatorname{Sym}^2P\to\Gamma^2P\).
On a diagonal square it is \(a^2\mapsto2\gamma_2(a)\);
on a mixed product it is \(ab\mapsto\gamma_1(a)\gamma_1(b)\).
No injectivity of this norm map is assumed; the
argument therefore applies also in characteristic two.

Both symmetric and divided squares preserve
projectives, by their free-module bases and preservation
of retracts. Dold--Kan shows that the simplicial terms
of \(P\) are projective, and normalization is a direct
summand of each term. Thus \(NR^2\) and \(NR^3\) are
degreewise projective. The Ext-acyclicity argument
used for (B-C2) gives
\[
 \operatorname{Ext}^1_{\mathcal S_k}
       (R^i,A^{2w,i-1})=0\qquad(i=2,3).
 \tag{D-30d}
\]
Indeed a degreewise-projective normalized source
has no positive abelian-category Ext into a horizontally
acyclic target, as proved there by degreewise splitting
of extensions and dimension shifting.

Apply \(\operatorname{Hom}_{\mathcal S_k}
(-,A^{2w,i-1})\) to
\(0\to C^i\to U^i\to R^i\to0\).
Equation (D-30d) extends each component of \(h\) to
\(\widetilde h^i:U^i\to A^{2w,i-1}\).
There is no quotient in degree four, so that component
needs no extension. Define
\[
                    g=d\widetilde h+\widetilde h d.
 \tag{D-30e}
\]
Then \(dg=gd\), and (D-30a) gives \(g|_{C(P)}=f\).
This is an actual simplicial linear extension, not
merely a derived morphism.

Every product of two arity-one elements lies in
\(C(P)\), as does every divided-power operation on
arity one whose output arity is two. Therefore
extending \(f\) on that subcomplex supplies exactly
all multiplicative and PD compatibility conditions
required through arity two. There are no additional
such conditions on a new arity-two element until
higher arity. This proves the asserted partial map.
Equivalently, it is a map of algebras after truncating
both source arity and the corresponding target weight
above two. The explicit terms (D-30b) show that no new
form-degree-one parameter has been introduced.
\end{proof}

\paragraph{Relation with the elementary integral test.}
For \(k=\Z\), \(P=K(\Z,2)\), the same extension also
follows without the acyclic-factorization hypothesis
on \(A^{w,1}\). In that case the norm in (D-30c) is
injective and \(R\simeq(P/2)[-3]\).
The sphere complex \(N(P/2)=(\Z/2)[2]\) has projective
dimension at most three in the abelian category of
nonnegative complexes: resolve the sphere by identity
disks and resolve \(\Z/2\) by two free abelian groups.
Thus \(\operatorname{Ext}^4_{\mathcal S_{\Z}}
(P/2,\underline\Z)=0\).
The bounded Hom comparison in degrees up to one,
using target rows zero through four, gives
\[
 \operatorname{Ext}^1_{\operatorname{Ch}(\mathcal S_{\Z})}
       (R,A^{2w,*})
 \cong H^1\operatorname{Tot}\operatorname{Hom}_{\mathcal S_{\Z}}
       (R,A^{2w,*})
 \cong\operatorname{Ext}^4_{\mathcal S_{\Z}}
       (P/2,\underline\Z)=0.
\]
The first equality uses the degreewise splitting of
extensions, justified by the same Ext vanishing as
(D-30d). Splitting the pushout extension of
\(0\to C(P)\to U\to R\to0\) along \(f\) gives
the desired strict extension. This alternate proof
is useful for checking the degree conventions;
the preceding proof gives the stated general-ring result.
